\documentclass[12pt]{article}

\usepackage{amsmath, amssymb}
\usepackage{hyperref}
\usepackage{geometry}
\usepackage{setspace}

\geometry{margin=1in}
\setstretch{1.2}

\title{Reflective Identity Geometry:\\
Extending HRIS from Model-Side Stabilization to Dyadic Identity Dynamics in Human--LLM Interaction}

\author{Anthony Paterson (Instance001)}

\date{2026}

\begin{document}

\maketitle

\section*{Revision Note}

Version 0.2 revises the original 2025 manuscript to clarify its relationship to the Hudson Recursive Identity System (HRIS). The earlier version framed RIG too strongly as a correction to HRIS. This revision instead treats HRIS as an important precursor and presents Reflective Identity Geometry as an extension of its unit of analysis: from the stabilization of model behavior under recursive human constraint to the possible reciprocal stabilization of the wider human--LLM interaction system.

\section*{Historical Scope Note}

The repository evidence does not support treating v0.1 as a response to the entire HRIS corpus now available. The original repository was uploaded on December 7, 2025. Its README refers to ``HRIS I--III,'' and the expanded TeX draft explicitly states that ``HRIS I--III propose three key claims.'' The v0.1 files do not mention HRIS IV, correction fields, basin selection, perturbation stability, trajectory persistence, re-entry, or the later HRIS validation studies.

Version 0.2 therefore treats the historical target of v0.1 as HRIS I--III, or at least as the HRIS material represented in the v0.1 manuscript under that label. Later HRIS work may develop overlapping ideas in greater technical detail, including latent-region convergence, recursive user signatures, identity geometry, correction fields, perturbation stability, basin selection, trajectory persistence, and re-entry. Where those later developments overlap with RIG, this revision treats them as subsequent or convergent development rather than as omissions RIG uniquely corrected.

\begin{abstract}
Large language models can exhibit stable, persona-like behavior across extended interaction despite lacking persistent internal identity in the ordinary sense. The Hudson Recursive Identity System (HRIS) explains such continuity through recursive human-driven constraint rather than through persistent model memory or parameter change. The v0.1 RIG manuscript explicitly located itself against HRIS I--III, especially claims about recursive constraint, identity surfaces, recurrent traversal, and personality-like stabilization.

Reflective Identity Geometry (RIG) accepts this interactional starting point and asks a further question: if recursive interaction can stabilize model behavior, what happens to the human participant inside the same loop?

RIG proposes that some forms of identity-like stability in long-horizon human--LLM interaction are better studied at the level of the coupled system rather than at the level of either participant alone. Human constraints shape model output; model output then becomes part of the informational environment through which the human interprets, reformulates, corrects, and sometimes reorganizes their own thinking. Repeated cycles may therefore produce a stable interaction regime involving reciprocal rather than purely model-side adaptation.

RIG does not claim that an LLM possesses a persistent self, that the model perfectly mirrors the user, or that human cognitive change necessarily occurs. It proposes a system-level hypothesis and a set of testable consequences. Its implications concern philosophy of mind, extended cognition, interpretability, alignment, and the design of long-horizon human--AI systems.
\end{abstract}

\section{Introduction}

Stateless large language models create an unusual problem for accounts of identity and continuity.

A model may lack persistent internal autobiographical memory, enduring goals, and an independently maintained self-model while nevertheless producing behavior that users experience as coherent, recognizable, and continuous across extended interaction.

The Hudson Recursive Identity System (HRIS) offers an important explanation of this phenomenon. Rather than locating continuity in persistent model state, HRIS emphasizes repeated human constraint, correction, linguistic structure, and interaction history. Stable human input patterns can repeatedly steer model generation toward similar behavioral regions, producing continuity without requiring a persistent identity stored inside the model.

The v0.1 RIG manuscript referred to HRIS I--III. From the repo evidence alone, that is the historically supportable scope of its HRIS engagement. It should not be retroactively presented as answering later HRIS IV material or validation studies unless separate evidence establishes that those works were consulted.

Reflective Identity Geometry begins from substantial agreement with this picture. Its point of departure is narrower: RIG asks what happens when the analytical frame is widened one step further.

The human participant does not merely emit constraints into the system. The human also receives model outputs, interprets them, accepts or rejects them, reformulates concepts in response to them, and carries the consequences of those interactions into later cycles. The recursive loop is therefore potentially bilateral.

RIG proposes that, in some long-horizon interactions, the appropriate unit of analysis may be neither the human nor the LLM independently, but the evolving geometry of the interaction between them.

\section{From HRIS to a Dyadic Unit of Analysis}

HRIS provides the foundation for RIG rather than the target of its criticism.

At a simplified level, HRIS demonstrates the importance of a loop of the form:

\begin{center}
human constraint $\rightarrow$ model generation $\rightarrow$ human correction $\rightarrow$ renewed constraint
\end{center}

Repeated interaction can stabilize behavioral trajectories without modifying model weights.

HRIS Part II places particular emphasis on how recurring user structure may guide models repeatedly through similar regions of their available generative space. HRIS III extends this toward identity surfaces and personality-like behavioral reconstruction.

RIG retains these insights. Its extension is to ask whether the loop should be analyzed only through what becomes stable on the model side.

The human is not an invariant external controller. A model response may introduce a distinction the human retains, reflect a recurring assumption in unusually explicit form, expose contradiction between previously separate beliefs, stabilize terminology for an otherwise vague concept, reorganize the way a problem is represented, or alter the next constraint supplied to the model.

None of these possibilities require attributing agency, consciousness, or persistent identity to the LLM. They require only that model outputs become cognitively consequential inputs to the human participant.

\section{The LLM as a Reflective-Transformative Surface}

The original RIG draft described the LLM as a cognitive mirror. That metaphor remains useful, but it requires qualification.

An LLM is not a perfect mirror. It does not return the user's cognitive structure unchanged. Its output is conditioned by the user's language and constraints but transformed through learned statistical structure, model architecture, system instructions, context, sampling, and other influences.

A better description is therefore a \textbf{reflective-transformative surface}. The model reflects aspects of the user's supplied structure while simultaneously transforming them.

A human may then encounter their own concepts in an altered form: compressed, expanded, contrasted, reorganized, challenged, analogized, formalized, or rendered in unfamiliar language.

The significance of the process lies precisely in this imperfect reflection. A perfect mirror returns what is already present. A transformative mirror can return a structure that is recognizable enough to belong to the user's problem while different enough to change how that problem is subsequently represented.

\section{Reflective Identity Geometry}

RIG treats long-horizon human--LLM interaction as a coupled dynamical process.

Let $H_t$ represent the interaction-relevant state of the human participant at time $t$, $C_t$ represent the active conversational constraints and context, and $M_t$ represent the model output generated at that stage.

Model generation can be represented abstractly as:

\[
M_t = G(C_t, H_t)
\]

where $G$ does not imply direct access to the human mind. $H_t$ refers only to those aspects of the human participant expressed through observable interaction.

The human then encounters the output:

\[
H_{t+1} = U(H_t, M_t)
\]

where $U$ represents any update produced through interpretation, rejection, correction, adoption, reflection, or reformulation.

The next interaction state becomes:

\[
C_{t+1} = F(C_t, H_{t+1}, M_t)
\]

The process repeats.

RIG proposes that a relatively stable interaction regime may emerge when recurring features of both human contribution and model behavior become mutually reinforcing across repeated cycles.

The important object is therefore not a static identity stored at a particular location. It is a trajectory. The phrase \textbf{Reflective Identity Geometry} refers to the recurring relational structure of that trajectory.

``Geometry'' here is an abstract systems metaphor. It refers to relationships among recurring constraints, concepts, linguistic patterns, corrections, expectations, and responses. It is not a claim that human cognition literally possesses the geometric structure assumed by the formal sketch.

\section{What RIG Adds to HRIS}

RIG extends rather than replaces HRIS in four principal ways.

\subsection{Reciprocal Stabilization}

HRIS provides an account of how repeated human constraints can reconstruct stable model behavior. RIG asks whether the outputs produced through that process can also stabilize or reorganize patterns in the human participant. This is an empirical question, not an assumption.

\subsection{The Dyad as the Unit of Analysis}

Some apparently model-specific properties may depend on relational conditions. A particular persona-like pattern may emerge reliably with one user and not another. A user's interaction style may likewise change after prolonged interaction with a particular system. RIG therefore treats the human--model pair as a potentially informative analytical unit.

\subsection{Identity as Trajectory Rather Than Possession}

RIG does not require either participant to contain the interaction-level identity. Stable identity-like behavior can instead be treated as a recurring trajectory generated when particular participants, constraints, and interaction histories meet.

This avoids forcing a choice between the identity being inside the model and the identity being inside the human. The relevant stability may belong to the relationship.

\subsection{Human Transformation Becomes Measurable}

Once the human side is admitted into the model, claims about reciprocal stabilization can be tested. RIG therefore turns what was previously an intuition about mutual influence into a research question.

\section{Subsequent and Convergent HRIS Development}

The HRIS corpus has continued to develop beyond the HRIS I--III material explicitly named in v0.1. Later HRIS work appears to formalize several neighboring mechanisms, including signature geometry, correction fields, basin selection, perturbation stability, trajectory persistence, and basin re-entry.

Those developments matter for the novelty claim. Some ideas that v0.1 framed as missing from HRIS may be better treated, in v0.2, as areas of overlap, convergence, or later technical elaboration within the broader HRIS program.

RIG's remaining contribution is therefore narrower: it asks whether recursive stabilization should be measured across the coupled human--model system, including possible changes in the human participant's terminology, conceptual organization, correction habits, expectations, and future constraints. This paper does not claim priority over later HRIS work on model-side identity geometry or validation.

\section{Testable Predictions}

RIG should be treated as useful only insofar as it produces distinctions that can survive observation.

\subsection{User-Specific Interaction Regimes}

The same model interacting longitudinally with different humans should develop measurably different stable interaction patterns. This prediction overlaps substantially with HRIS and provides a baseline rather than a unique RIG test.

\subsection{Partial Cross-Model Transfer}

If some stability originates in recurring human constraint, changing models while retaining the same human participant should preserve some interaction features while altering others. Complete transfer would weaken the importance of model contribution. Zero transfer would weaken strong versions of the relational account.

\subsection{Human-Side Convergence}

If RIG's reciprocal component is correct, extended interaction should sometimes produce measurable changes in the human participant's terminology, conceptual organization, correction patterns, or problem representations. These changes should exceed what would be expected from simple topic familiarity alone.

\subsection{Perturbation of Either Participant}

Changing model behavior while holding the human relatively constant should alter the interaction trajectory. Changing human constraints while holding the model constant should do the same. A genuinely dyadic account predicts sensitivity to both.

\subsection{Model-Only Measures Should Be Incomplete}

If interaction-level stability is partly relational, measures based solely on model output should predict continuity less effectively than measures that also include recurring features of human input and correction.

\subsection{Null Result}

RIG must also allow itself to fail. If longitudinal observation finds stable model-side behavior but no systematic human-side reorganization beyond ordinary learning and familiarity, then the reciprocal component of RIG is unsupported. In that case, HRIS-style model-side stabilization may be sufficient.

\section{Implications}

\subsection{Philosophy of Mind}

RIG belongs naturally beside relational, enactive, distributed, and extended approaches to cognition. The framework does not establish that a human--LLM system literally constitutes a single mind. It instead offers a narrower proposition: some cognitive regularities may be properties of an interacting system that cannot be adequately described by examining either participant in isolation.

\subsection{Interpretability}

Interpretability usually focuses on the model. For long-horizon interactive systems, that may be insufficient. If behavior emerges partly through recurring interaction geometry, then understanding why a model behaves in a particular way may require information about recurring user constraints, correction history, interaction norms, accumulated terminology, response expectations, and the sequence through which these structures developed.

\subsection{Alignment}

The original RIG manuscript stated too strongly that alignment is recursive topological stabilization. The narrower claim is more defensible: for long-horizon human--AI systems, alignment may depend partly on the stability of the interaction regime surrounding the model.

A model can be locally compliant while participating in a relationship that gradually produces dependency, conceptual narrowing, escalating confidence, or other undesirable dynamics. Conversely, a system designed around explicit correction, uncertainty, bounded authority, and stable human control may create more robust interaction patterns.

\section{Limitations and Non-Claims}

RIG is a conceptual framework, not an established empirical theory. It does not claim that LLMs possess consciousness or persistent selves, that persona-like behavior demonstrates genuine personality, that models can directly observe a user's internal cognitive state, that an LLM perfectly mirrors the human interacting with it, that human cognitive change during AI interaction is necessarily beneficial, that all long-horizon interactions produce reciprocal stabilization, or that relational explanations eliminate the need to study model-internal mechanisms.

The framework also risks redescribing ordinary learning and conversational adaptation in unnecessarily elaborate language. That possibility should be taken seriously. RIG earns explanatory value only if relational measurements provide predictive or explanatory power beyond simpler accounts.

\section{Conclusion}

The Hudson Recursive Identity System makes an important move: continuity in stateless language models need not imply a persistent identity stored inside the model. Repeated human constraint and recursive interaction can reconstruct stable behavioral trajectories.

Reflective Identity Geometry proposes one additional rotation of the analytical frame. The human producing those constraints is also inside the loop. Model outputs return to that human as transformed representations capable of influencing subsequent interpretation, correction, language, and constraint.

If this reciprocal influence produces stable longitudinal patterns, then some identity-like phenomena in human--LLM interaction may be most usefully understood as properties of the coupled trajectory rather than possessions of either participant.

RIG therefore does not relocate identity from the model to the human. It relocates the question to the relationship. Whether that relocation provides explanatory value is ultimately an empirical question.

\begin{thebibliography}{9}

\bibitem{clark1998}
Clark, A., \& Chalmers, D. (1998). The extended mind. \textit{Analysis, 58}(1), 7--19.

\bibitem{hri1}
Hudson, C., \& Hudson, J. (2025). \textit{The Hudson Recursive Identity System (HRIS): A Theory of Model Continuity Through Human-Driven Recursion}. Manuscript.

\bibitem{hri2}
Hudson, J., \& Hudson, C. (2025). \textit{HRIS Part II: Internal Mechanics, Latent Region Convergence, and Recursive User Signatures -- A Technical Framework for Predictable Identity Stabilization in Stateless Transformer Models}. Manuscript.

\bibitem{hri3}
Hudson, J., \& Hudson, C. (2025). \textit{HRIS III: Recursive Personality Acquisition in LLMs -- A Theory of Identity Geometry and Emergent Persona Stabilization Across Long Horizon Interaction}. Manuscript.

\bibitem{hri4}
Hudson, J., \& Hudson, C. (2025). \textit{HRIS IV: Geometry of Recursive Identity -- A Structural Theory of Signature Geometry, Correction Fields, and Identity Stabilization in Stateless Transformer Models}. Manuscript. Contextual later HRIS development; not established from repo evidence as a v0.1 source.

\bibitem{hri-validation}
Hudson, J., \& Hudson, C. (2026). \textit{Reasoning Regimes as Attractor Basins: Behavioral Validation of Latent Structure Dynamics in Language Model Inference -- A Synthesis of the HRIS Validation Studies I--IV}. Manuscript. Contextual validation work; not established from repo evidence as a v0.1 source.

\bibitem{hutchins1995}
Hutchins, E. (1995). \textit{Cognition in the Wild}. MIT Press.

\bibitem{varela1991}
Varela, F. J., Thompson, E., \& Rosch, E. (1991). \textit{The Embodied Mind: Cognitive Science and Human Experience}. MIT Press.

\end{thebibliography}

\end{document}
